\section{Reconstructing the Published-Theorem Baseline}
\label{sec:baseline}

The phrase ``published baseline'' can obscure two logically distinct objects.
The first is the high-order commutator theorem published by Childs et al.
\cite{childs2021trotter}; the second is our fixed, machine-evaluated
instantiation of that theorem.  The number 393 is not a table entry quoted
from the 2021 paper.  It is obtained by applying the published framework to
\cref{eq:hamiltonian,eq:benchmark,eq:suzuki4} with a specified norming and
interval procedure.  We call it the \emph{pinned control}.

\subsection{Expansion-first local norming}

The high-order theorem represents product-formula error by partial sums and
nested commutators.  For the pinned control, every partial matching sum is
expanded into local Heisenberg bonds, zero commutators are removed, and each
surviving concrete Pauli expansion is bounded by its Pauli-$\ell_1$ norm.
The resulting site-density upper bound is stored as an exact rational number
whose outward decimal is
\begin{equation}
 \alpha_{\mathrm{base}}\le164.9759169527440.
\end{equation}
The calculation contains 8,316 theorem terms and 1,024 expanded commutator
keys.  All algebraic product-formula coefficients are enclosed by rational
intervals; no rounded decimal is used as the source of the result.

The theorem admits a discrete family of center choices controlling how the
high-order integral representation is partitioned.  We evaluate all 31
admissible centers.  Center 20 gives the smallest certified control, and
inverting its global bound at tolerance $10^{-6}$ yields
\begin{equation}
 r_{\mathrm{base}}=393,\qquad G(r_{\mathrm{base}})=11{,}791.
\end{equation}
The main certificate stores both the exact site-density rational and the
chosen center.  A deep verification mode regenerates the center scan rather
than trusting these stored labels.

\subsection{Where the baseline loses information}

The baseline is rigorous, general, and already exploits exact commutation.
Its looseness in this instance arises from the order of operations.  Once an
abstract contribution is replaced by an independent positive norm, later
knowledge cannot recover cancellation between that contribution and another
one.  Four losses are particularly important:

\begin{enumerate}
\item distinct free words can map to the same nested-commutator or Pauli
  contribution with opposite signs;
\item distinct translated local clusters can share a canonical translation
  representative;
\item different commutator paths can produce exactly the same Pauli string;
\item different surviving Pauli strings can anticommute, causing cross terms
  to vanish when the group norm is squared.
\end{enumerate}

The proposed compiler postpones the irreversible inequality until these
relations have been exposed.  Importantly, it does not treat the baseline as
incorrect.  It constructs a stronger proof for the same circuit and then
verifies both sides under a common benchmark and cost model.

\begin{table}[t]
\centering
\caption{Logical separation between the published theorem and the pinned
finite-instance control.}
\label{tab:baseline-separation}
\begin{tabularx}{\textwidth}{@{}lXX@{}}
\toprule
 & Published theory & Pinned control used here \\
\midrule
Source & Childs et al. commutator-scaling framework & Exact interval
instantiation for the fixed Heisenberg benchmark \\
Scope & General Hamiltonian fragments and formula order & $L=12$, four
matchings, five-copy $S_4$, $T=1$, $\eps=10^{-6}$ \\
Norming & Theorem permits commutator-sensitive evaluation & Expanded local
Pauli-$\ell_1$ procedure fixed by the certificate \\
Finite number & No published 393-step table entry & 31 centers scanned;
center 20 yields 393 steps \\
Role & Mathematical foundation & Auditable comparison control \\
\bottomrule
\end{tabularx}
\end{table}
